% !TITLE 上李邕
% !AUTHOR 李白
% !DYNASTY 唐
% !ORDER 4

\begin{poem}
大鹏\textsuperscript{1}一日同风起，扶摇\textsuperscript{2}直上九万里。
假令风歇时下来，犹能簸却\textsuperscript{3}沧溟\textsuperscript{4}水。
世人见我恒殊调\textsuperscript{5}，闻余大言皆冷笑。
宣父\textsuperscript{6}犹能畏后生，丈夫未可轻年少。
\end{poem}

\begin{annotations}
\notetext{1}{大鹏：传说中的巨鸟，出自《庄子·逍遥游》。大鹏的翅膀像天边的云，乘风飞上九万里高空。李白一生最爱以大鹏自喻。}
\notetext{2}{扶摇：盘旋而上的暴风。大鹏乘着扶摇旋风直上高空。}
\notetext{3}{簸却：摇动、掀翻。即便风停落下来，大鹏也能搅动大海的水。簸（bǒ），颠簸摇动。}
\notetext{4}{沧溟：大海。}
\notetext{5}{殊调：与众不同的格调、言论。世人看我总是格调怪异，听到我的豪言壮语都冷笑。}
\notetext{6}{宣父：指孔子。孔子曾说过“后生可畏”——年少的人值得敬畏。李白抬出孔子来为自己辩护，说大丈夫不该轻视少年。}
\end{annotations}

\begin{appreciation}
这是李白青年时代的作品，写于他出蜀漫游、求谒名士李邕之时。那一年他大约二十五岁，羽翼初丰，志气正盛。

李邕是盛唐著名的书法家和文学家，官至北海太守，以爱才养士闻名。年轻的李白慕名前去拜访，却因言辞放诞、不拘礼法而惹恼了这位前辈。李邕拂袖而去，李白则写下了这首诗作为回应——不是道歉，而是宣言。

大鹏一日同风起，扶摇直上九万里——开篇就是《庄子·逍遥游》中的意象。大鹏是李白一生最钟爱的自喻，从青年到暮年，这个意象贯穿了他的全部诗作。但值得注意的是，青年李白笔下的大鹏，是一种确定无疑的飞翔：风至则起，直上九万里，没有任何犹豫。这种自信近乎狂傲，却也是盛唐气象在个体身上的投射：一个年轻人相信，只要时机到来，他必然能够一飞冲天。

假令风歇时下来，犹能簸却沧溟水——即便风停了，大鹏落下来，也能搅动沧海之水。这是何等的自我期许。李白不是在说自己一定能成功，他是在说：即便失败，我的失败也必将惊天动地。这种姿态，后来的人学不来，因为那是一种时代赋予的底气。

世人见我恒殊调，闻余大言皆冷笑——世人看我总是格调奇异，听到我的豪言壮语都冷笑。宣父犹能畏后生，丈夫未可轻年少——孔子尚且敬畏年轻人，大丈夫不可轻视少年。宣父是孔子的尊称。李白抬出孔子来压李邕，这不是狡辩，是一种骨子里的不服输：你可以不欣赏我，但你不能否定我的未来。

这首诗里的李白，和后来我们所熟知的那个醉卧长安、放逐夜郎的李白，是同一个人，又不是同一个人。他还没有经历人生的跌宕，还没有写出天生我材必有用时的那种在绝望中硬撑的豪迈。这里的大鹏，是未经风雨的大鹏，是确信自己必将高飞的大鹏。这种确信本身，就是青春最珍贵的东西。
\end{appreciation}
